\chapter{2005年全国II卷（文）}

\section{单选题}

\begin{problem}
函数 \( f(x) = |\sin x + \cos x| \) 的最小正周期是（\quad）

\choices
    {\(\frac{\pi}{4}\)}
    {\(\frac{\pi}{2}\)}
    {\(\pi\)}
    {\(2\pi\)}
\end{problem}

\begin{answer}C\end{answer}

\begin{solution}
利用辅助角公式，\(\sin x+\cos x=\sqrt{2}\sin\left(x+\frac{\pi}{4}\right)\)，所以 \(f(x)=\sqrt{2}\left|\sin\left(x+\frac{\pi}{4}\right)\right|\). 函数 \(|\sin t|\) 的零点为 \(k\pi\)，相邻零点间距为 \(\pi\)，且 \(|\sin(t+\pi)|=|\sin t|\)，故其最小正周期是 \(\pi\). 平移自变量和乘以非零常数都不改变周期，因此 \(f(x)\) 的最小正周期为 \(\pi\)，故选 C.
\end{solution}

\begin{problem}
正方体 \(ABCD-A_{1}B_{1}C_{1}D_{1}\) 中，\(P\)，\(Q\)，\(R\) 分别是 \(AB\)，\(AD\)，\(B_{1}C_{1}\) 的中点. 那么，正方体的过 \(P\)，\(Q\)，\(R\) 的截面图形是（\quad）

\choices
    {三角形}
    {四边形}
    {五边形}
    {六边形}
\end{problem}

\begin{answer}D\end{answer}

\begin{solution}
设正方体边长为 \(1\)，取 \(A=(0,0,0)\)，\(B=(1,0,0)\)，\(D=(0,1,0)\)，\(A_{1}=(0,0,1)\). 则 \(P=\left(\frac{1}{2},0,0\right)\)，\(Q=\left(0,\frac{1}{2},0\right)\)，\(R=\left(1,\frac{1}{2},1\right)\). 由三点式可得过三点的平面方程为 \(x+y-z=\frac{1}{2}\). 该平面与正方体的棱分别交于底面上的 \(P\)、\(Q\)，棱 \(BB_{1}\) 上的一点，棱 \(B_{1}C_{1}\) 上的点 \(R\)，棱 \(C_{1}D_{1}\) 上的一点以及棱 \(DD_{1}\) 上的一点. 这六个交点依次连接构成截面，且没有其他交点，所以截面是六边形，故选 D.
\end{solution}

\begin{problem}
函数 \( y = x^{2} - 1 \)（\(x \leqslant 0\)）的反函数是（\quad）

\choices
    {\( y = \sqrt{x + 1} \)（\(x \geqslant -1\)）}
    {\( y = -\sqrt{x + 1} \)（\(x \geqslant -1\)）}
    {\( y = \sqrt{x + 1} \)（\(x \geqslant 0\)）}
    {\( y = -\sqrt{x + 1} \)（\(x \geqslant 0\)）}
\end{problem}

\begin{answer}B\end{answer}

\begin{solution}
原函数的定义域为 \((-\infty,0]\)，在该区间上 \(y=x^{2}-1\) 单调递减，值域为 \([-1,+\infty)\)，因而存在反函数. 由 \(y=x^{2}-1\) 得 \(x^{2}=y+1\). 原函数的自变量满足 \(x\leqslant0\)，所以应取 \(x=-\sqrt{y+1}\)，同时要求 \(y\geqslant-1\). 交换自变量与因变量，反函数为 \(y=-\sqrt{x+1}\)，定义域为 \(x\geqslant-1\)，故选 B.
\end{solution}

\begin{problem}
已知函数 \( y = \tan \omega x \) 在 \(\left(-\frac{\pi}{2}, \frac{\pi}{2}\right)\) 内是减函数，则（\quad）

\choices
    {\(0 < \omega \leqslant 1\)}
    {\(-1 \leqslant \omega < 0\)}
    {\(\omega \geqslant 1\)}
    {\(\omega \leqslant -1\)}
\end{problem}

\begin{answer}B\end{answer}

\begin{solution}
要使函数在整个开区间内有定义，必须没有点满足 \(\omega x=\frac{\pi}{2}+k\pi\). 若 \(|\omega|>1\)，则 \(x=\frac{\pi}{2|\omega|}\) 或其相反数落在 \(\left(-\frac{\pi}{2},\frac{\pi}{2}\right)\) 内，并使正切无定义；若 \(|\omega|\leqslant1\)，则对区间内任意 \(x\)，有 \(|\omega x|<\frac{\pi}{2}\)，函数有定义. 因此必须有 \(|\omega|\leqslant1\). 又 \(y'=\omega\sec^{2}(\omega x)\)，且 \(\sec^{2}(\omega x)>0\)，所以函数为减函数当且仅当 \(\omega<0\). 综合得 \(-1\leqslant\omega<0\)，故选 B.
\end{solution}

\begin{problem}
抛物线 \(x^{2} = 4y\) 上一点 \(A\) 的纵坐标为 \(4\)，则点 \(A\) 与抛物线焦点的距离为（\quad）

\choices
    {\(2\)}
    {\(3\)}
    {\(4\)}
    {\(5\)}
\end{problem}

\begin{answer}D\end{answer}

\begin{solution}
将抛物线写成标准形式 \(x^{2}=4ay\)，可知 \(a=1\)，所以焦点为 \(F=(0,1)\). 点 \(A\) 的纵坐标为 \(4\)，代入方程得 \(x^{2}=16\)，故 \(A=(4,4)\) 或 \(A=(-4,4)\). 两种情况下 \(AF=\sqrt{(\pm4-0)^{2}+(4-1)^{2}}=\sqrt{16+9}=5\)，故选 D.
\end{solution}

\begin{problem}
双曲线 \(\frac{x^2}{4} - \frac{y^2}{9} = 1\) 的渐近线方程是（\quad）

\choices
    {\( y = \pm \frac{2}{3} x \)}
    {\( y = \pm \frac{4}{9} x \)}
    {\( y = \pm \frac{3}{2} x \)}
    {\(\sqrt{3}\)}
\end{problem}

\begin{answer}C\end{answer}

\begin{solution}
双曲线 \(\frac{x^{2}}{a^{2}}-\frac{y^{2}}{b^{2}}=1\) 的渐近线由最高次项方程 \(\frac{x^{2}}{a^{2}}-\frac{y^{2}}{b^{2}}=0\) 得出，即 \(y=\pm\frac{b}{a}x\). 本题中 \(a=2\)，\(b=3\)，所以渐近线为 \(y=\pm\frac{3}{2}x\)，故选 C.
\end{solution}

\begin{problem}
如果数列 \(\{a_{n}\}\) 是等差数列，则（\quad）

\choices
    {\(a_1 + a_8 < a_4 + a_5\)}
    {\(a_1 + a_8 = a_4 + a_5\)}
    {\(a_{1} + a_{8} > a_{4} + a_{5}\)}
    {\(a_{1}a_{8} = a_{4}a_{5}\)}
\end{problem}

\begin{answer}B\end{answer}

\begin{solution}
设等差数列的首项为 \(a_{1}\)，公差为 \(d\)，则 \(a_{n}=a_{1}+(n-1)d\). 因此 \(a_{8}=a_{1}+7d\)，从而 \(a_{1}+a_{8}=2a_{1}+7d\). 另一方面，\(a_{4}+a_{5}=\left(a_{1}+3d\right)+\left(a_{1}+4d\right)=2a_{1}+7d\). 两式相等，所以 \(a_{1}+a_{8}=a_{4}+a_{5}\)，故选 B.
\end{solution}

\begin{problem}
\((x - \sqrt{2} y)^{10}\) 的展开式中 \(x^{6}y^{4}\) 项的系数是（\quad）

\choices
    {\(840\)}
    {\(-840\)}
    {\(210\)}
    {\(-210\)}
\end{problem}

\begin{answer}A\end{answer}

\begin{solution}
二项式展开的通项为 \(\mathrm C_{10}^{k}x^{10-k}\left(-\sqrt{2}y\right)^{k}\). 要得到 \(x^{6}y^{4}\)，必须有 \(10-k=6\)，所以 \(k=4\). 此时对应项的系数为 \(\mathrm C_{10}^{4}\left(-\sqrt{2}\right)^{4}=210\times4=840\)，故选 A.
\end{solution}

\begin{problem}
已知点 \(A(\sqrt{3}, 1)\)，\(B(0, 0)\)，\(C(\sqrt{3}, 0)\). 设 \(\angle BAC\) 的平分线 \(AE\) 与 \(BC\) 相交于 \(E\)，那么有 \(\overrightarrow{BC} = \lambda \overrightarrow{CE}\)，其中 \(\lambda\) 等于（\quad）

\choices
    {\(2\)}
    {\(\frac{1}{2}\)}
    {\(-3\)}
    {\(-\frac{1}{3}\)}
\end{problem}

\begin{answer}C\end{answer}

\begin{solution}
由坐标距离公式，\(AB=\sqrt{(\sqrt3)^{2}+1^{2}}=2\)，\(AC=1\)，\(BC=\sqrt3\). 由角平分线定理，\(\frac{BE}{EC}=\frac{AB}{AC}=2\). 设 \(EC=t\)，则 \(BE=2t\)，而 \(BE+EC=BC\)，所以 \(3t=\sqrt3\)，即 \(EC=\frac{\sqrt3}{3}\). 点 \(E\) 在线段 \(BC\) 内部，因此 \(\overrightarrow{CE}\) 与 \(\overrightarrow{BC}\) 方向相反；又 \(|\overrightarrow{BC}|=\sqrt3=3|\overrightarrow{CE}|\). 所以 \(\overrightarrow{BC}=-3\overrightarrow{CE}\)，即 \(\lambda=-3\)，故选 C.
\end{solution}

\begin{problem}
已知集合 \(M = \{x \mid x^2 - 3x - 28 \leqslant 0\}\)，\(N = \{x \mid x^2 - x - 6 > 0\}\)，则 \(M \cap N\) 为（\quad）

\choices
    {\(\left\{x\mid -4\leqslant x < -2\text{\ 或\ }3 < x\leqslant 7\right\}\)}
    {\(\left\{x\mid -4 < x\leqslant -2\text{\ 或\ }3\leqslant x < 7\right\}\)}
    {\(\left\{x\mid x\leqslant -2\text{\ 或\ }x > 3\right\}\)}
    {\(\left\{x\mid x < -2\text{\ 或\ }x\geqslant 3\right\}\)}
\end{problem}

\begin{answer}A\end{answer}

\begin{solution}
因 \(x^{2}-3x-28=(x-7)(x+4)\)，且二次项系数为正，所以 \(x^{2}-3x-28\leqslant0\) 的解集为 \(M=[-4,7]\). 又因 \(x^{2}-x-6=(x-3)(x+2)\)，所以 \(x^{2}-x-6>0\) 的解集为 \(N=(-\infty,-2)\cup(3,+\infty)\)，端点 \(-2\)、\(3\) 因严格不等式不取. 取交集得 \(M\cap N=[-4,-2)\cup(3,7]\)，故选 A.
\end{solution}

\begin{problem}
点 \(P\) 在平面上作匀速直线运动. 速度向量 \(\bs{v} = (4, -3)\)（即点 \(P\) 的运动方向与 \(\bs{v}\) 相同，且每秒移动的距离为 \(|\bs{v}|\text{ 个单位}\)）. 设开始时点 \(P\) 的坐标为 \((-10, 10)\)，则 \(5\text{ 秒}\) 后点 \(P\) 的坐标为（\quad）

\choices
    {\((-2, 4)\)}
    {\((-30, 25)\)}
    {\((10, -5)\)}
    {\((5, -10)\)}
\end{problem}

\begin{answer}C\end{answer}

\begin{solution}
速度向量的模为 \(|\bs{v}|=\sqrt{4^{2}+(-3)^{2}}=5\). 题设说明每秒移动的距离也是 \(5\)，并且运动方向与 \(\bs{v}\) 相同，因此每秒的位移向量就是 \(\bs{v}=(4,-3)\). 经过 \(5\) 秒，位移为 \(5\bs{v}=(20,-15)\)，所以终点坐标为 \((-10,10)+(20,-15)=(10,-5)\)，故选 C.
\end{solution}

\begin{problem}
\(\triangle ABC\) 的顶点 \(B\) 在平面 \(\alpha\) 内，\(A\)，\(C\) 在 \(\alpha\) 的同一侧，\(AB\)，\(BC\) 与 \(\alpha\) 所成的角分别是 \(30^{\circ}\) 和 \(45^{\circ}\). 若 \(AB = 3\)，\(BC = 4\sqrt{2}\)，\(AC = 5\)，则 \(AC\) 与 \(\alpha\) 所成的角为（\quad）

\choices
    {\(60^{\circ}\)}
    {\(45^{\circ}\)}
    {\(30^{\circ}\)}
    {\(15^{\circ}\)}
\end{problem}

\begin{answer}C\end{answer}

\begin{solution}
设 \(A\)、\(C\) 到平面 \(\alpha\) 的距离分别为 \(h_A\)、\(h_C\). 由于 \(B\in\alpha\)，由线面角的定义，\(h_A=AB\sin30^{\circ}=3\times\frac12=\frac32\)，\(h_C=BC\sin45^{\circ}=4\sqrt2\times\frac{\sqrt2}{2}=4\). 因为 \(A\)、\(C\) 在平面同侧，线段 \(AC\) 的垂直分量为两点到平面的距离之差，即 \(h_C-h_A=\frac52\). 设 \(AC\) 与平面 \(\alpha\) 所成的角为 \(\theta\)，则 \(AC\sin\theta=\frac52\). 代入 \(AC=5\)，得 \(\sin\theta=\frac12\). 线面角取 \([0^{\circ},90^{\circ}]\) 内的角，所以 \(\theta=30^{\circ}\)，故选 C.
\end{solution}

\section{填空题}

\begin{problem}
在 \(\frac{8}{3}\) 和 \(\frac{27}{2}\) 之间插入三个数，使这五个数成等比数列，则插入的三个数的乘积为\fillinblank{}.
\end{problem}

\begin{answer}\(216\)\end{answer}

\begin{solution}
设五个数依次为 \(a_{1}\)，\(a_{2}\)，\(a_{3}\)，\(a_{4}\)，\(a_{5}\)，公比为 \(q\). 则 \(a_{1}=\frac83\)，\(a_{5}=\frac{27}{2}\)，从而 \(a_{5}=a_{1}q^{4}\)，即 \(q^{4}=\frac{27/2}{8/3}=\frac{81}{16}\)，所以 \(q=\frac32\) 或 \(q=-\frac32\). 无论公比取哪一个值，中项性质都给出 \(a_{3}^{2}=a_{1}a_{5}=\frac83\cdot\frac{27}{2}=36\). 又 \(a_{3}=a_{1}q^{2}>0\)，故 \(a_{3}=6\). 因为 \(a_{2}a_{4}=a_{3}^{2}\)，所以插入的三个数的乘积为 \(a_{2}a_{3}a_{4}=a_{3}^{3}=6^{3}=216\).
\end{solution}

\begin{problem}
圆心为 \((1,2)\) 且与直线 \(5x - 12y - 7 = 0\) 相切的圆的方程为\fillinblank{}.
\end{problem}

\begin{answer}\(\left(x-1\right)^{2}+\left(y-2\right)^{2}=4\)\end{answer}

\begin{solution}
圆与直线相切时，圆心到直线的距离就是半径. 圆心 \(O=(1,2)\) 到直线 \(5x-12y-7=0\) 的距离为 \(r=\frac{|5\times1-12\times2-7|}{\sqrt{5^{2}+(-12)^{2}}}=\frac{26}{13}=2\). 因此圆心为 \((1,2)\)、半径为 \(2\) 的圆方程为 \(\left(x-1\right)^{2}+\left(y-2\right)^{2}=2^{2}=4\).
\end{solution}

\begin{problem}
在由数字 \(0\)，\(1\)，\(2\)，\(3\)，\(4\)，\(5\) 所组成的没有重复数字的四位数中，不能被 \(5\) 整除的数共有\fillinblank{}\(\text{ 个}\).
\end{problem}

\begin{answer}\(192\)\end{answer}

\begin{solution}
四位数不能被 \(5\) 整除，末位不能为 \(0\) 或 \(5\)，所以末位有 \(4\) 种选法. 固定末位后，剩余五个数字中有四个非零数字，首位有 \(4\) 种选法；首位确定后，剩余四个数字中任选两个依次放在中间两位，有 \(4\times3\) 种选法. 根据乘法原理，所求个数为 \(4\times4\times4\times3=192\).
\end{solution}

\begin{problem}
下面是关于三棱锥的四个命题：

\begin{circlelist}
\item 底面是等边三角形，侧面与底面所成的二面角都相等的三棱锥是正三棱锥；
\item 底面是等边三角形，侧面都是等腰三角形的三棱锥是正三棱锥；
\item 底面是等边三角形，侧面的面积都相等的三棱锥是正三棱锥；
\item 侧棱与底面所成的角都相等，且侧面与底面所成的二面角都相等的三棱锥是正三棱锥.
\end{circlelist}

其中，真命题的编号是\fillinblank{}. （写出所有真命题的编号）
\end{problem}

\begin{answer}\(\circled{1}\circled{4}\)\end{answer}

\begin{solution}
设三棱锥顶点为 \(S\)，底面为等边三角形 \(ABC\)，点 \(H\) 为 \(S\) 在底面上的射影，记 \(SH=h>0\).

对于命题 \(\circled{1}\)，在过 \(SH\) 且垂直于某一条底边的截面内，二面角的平面角由底边内侧方向和截面内的侧棱方向构成. 若 \(H\) 在该底边所对应的内侧半平面内，则该角的正切为 \(h\) 除以 \(H\) 到该底边所在直线的距离；若 \(H\) 在外侧，则相应的内二面角为钝角. 三个二面角相等，因而 \(H\) 必须同时处在三条边的同侧，并且到三条边所在直线的距离相等. 三个外侧半平面不可能有公共点，所以 \(H\) 只能在三角形 \(ABC\) 内，且到三边距离相等，即 \(H\) 是内心. 等边三角形的内心也是外心，所以 \(HA=HB=HC\). 于是 \(SA^{2}=SH^{2}+HA^{2}\)，\(SB^{2}=SH^{2}+HB^{2}\)，\(SC^{2}=SH^{2}+HC^{2}\)，从而 \(SA=SB=SC\)，三棱锥为正三棱锥，命题 \(\circled{1}\) 正确.

命题 \(\circled{2}\) 不正确. 取底面 \(A=(-1,0,0)\)，\(B=(1,0,0)\)，\(C=(0,\sqrt3,0)\)，顶点 \(S=(0,0,1)\). 底面三边长均为 \(2\)，且 \(SA=SB=\sqrt2\)，\(SC=2\)，\(BC=CA=2\). 因此 \(\triangle SAB\) 中 \(SA=SB\)，\(\triangle SBC\) 中 \(SC=BC\)，\(\triangle SCA\) 中 \(SC=CA\)，三个侧面都是等腰三角形；但是 \(SA\ne SC\)，且顶点在底面的射影 \((0,0,0)\) 不是等边三角形的中心，所以该三棱锥不是正三棱锥.

命题 \(\circled{3}\) 不正确. 仍取上述底面，令 \(H=(0,-\sqrt3,0)\)，\(S=(0,-\sqrt3,1)\). 点 \(H\) 到三条边所在直线的距离都为 \(\sqrt3\)：到 \(AB\) 的距离显然为 \(\sqrt3\)，到 \(BC\)、\(CA\) 的距离由点到直线距离公式也均为 \(\sqrt3\). 因为 \(SH=1\)，所以点 \(S\) 到三条底边所在直线的距离都为 \(\sqrt{(\sqrt3)^{2}+1^{2}}=2\). 三个侧面分别以底边为底、以这些距离为高，且底边长均为 \(2\)，所以三个侧面面积均为 \(\frac12\times2\times2=2\). 但 \(H\) 是底面三角形的旁心而不是中心，故此三棱锥不是正三棱锥，命题 \(\circled{3}\) 错误.

对于命题 \(\circled{4}\)，若三条侧棱与底面所成的角相等，则由线面角定义有 \(\sin\theta=\frac{SH}{SA}=\frac{SH}{SB}=\frac{SH}{SC}\). 因为 \(SH>0\)，所以 \(SA=SB=SC\). 在直角三角形 \(SHA\)、\(SHB\)、\(SHC\) 中，得 \(HA=HB=HC\)，所以 \(H\) 是底面等边三角形的外心，也就是其中心. 底面为等边三角形且顶点射影为底面中心，故三棱锥为正三棱锥，命题 \(\circled{4}\) 正确. 因此真命题的编号为 \(\circled{1}\)、\(\circled{4}\).
\end{solution}

\section{解答题}

\begin{problem}
已知 \(\alpha\) 为第二象限的角，\(\sin \alpha = \frac{3}{5}\)，\(\beta\) 为第一象限的角，\(\cos \beta = \frac{5}{13}\). 求 \(\tan (2\alpha - \beta)\) 的值.
\end{problem}

\begin{solution}
因为 \(\alpha\) 在第二象限，\(\cos\alpha<0\)，由 \(\sin^{2}\alpha+\cos^{2}\alpha=1\) 得 \(\cos\alpha=-\frac45\)，从而 \(\tan\alpha=-\frac34\). 因此
\(\tan2\alpha=\frac{2\tan\alpha}{1-\tan^{2}\alpha}=\frac{-\frac32}{1-\frac{9}{16}}=-\frac{24}{7}\).

因为 \(\beta\) 在第一象限，\(\sin\beta>0\)，所以 \(\sin\beta=\sqrt{1-\frac{25}{169}}=\frac{12}{13}\)，从而 \(\tan\beta=\frac{12}{5}\). 和差角正切公式的分母为 \(1+\tan2\alpha\tan\beta=1-\frac{288}{35}=-\frac{253}{35}\ne0\)，故
\(\tan(2\alpha-\beta)=\frac{\tan2\alpha-\tan\beta}{1+\tan2\alpha\tan\beta}=\frac{-\frac{24}{7}-\frac{12}{5}}{1-\frac{24}{7}\cdot\frac{12}{5}}=\frac{-\frac{204}{35}}{-\frac{253}{35}}=\frac{204}{253}\).
\end{solution}

\begin{problem}
甲，乙两队进行一场排球比赛，根据以往经验，单局比赛甲队胜乙队的概率为 \(0.6\)，本场比赛采用五局三胜制，即先胜三局的队获胜，比赛结束，设各局比赛相互间没有影响. 求：

\begin{enumerate}
\item 前三局比赛甲队领先的概率；
\item 本场比赛乙队以 \(3:2\) 取胜的概率. （精确到 \(0.001\)）
\end{enumerate}
\end{problem}

\begin{solution}
设每局甲队获胜事件的概率为 \(p=0.6\)，乙队获胜事件的概率为 \(q=1-p=0.4\). 各局相互独立.

\begin{enumerate}
\item 前三局后甲队领先，表示前三局至少胜两局，包含恰胜两局和三局全胜两种互斥情况. 因此所求概率为 \(\mathrm C_{3}^{2}p^{2}q+p^{3}=3\times0.6^{2}\times0.4+0.6^{3}=0.432+0.216=0.648\).
\item 乙队以 \(3:2\) 取胜，比赛必须进行五局，且前四局乙队恰胜两局、甲队恰胜两局，第五局乙队获胜. 前四局中乙队胜两局的选法有 \(\mathrm C_{4}^{2}\) 种，所以所求概率为 \(\mathrm C_{4}^{2}q^{2}p^{2}q=6\times0.4^{2}\times0.6^{2}\times0.4=0.13824\). 精确到 \(0.001\) 得 \(0.138\).
\end{enumerate}
\end{solution}

\begin{problem}
已知 \(\left\{a_{n}\right\}\) 是各项均为正数的等差数列，\(\lg a_{1}\)，\(\lg a_{2}\)，\(\lg a_{4}\) 成等差数列. 又 \(b_{n} = \frac{1}{a_{2^{n}}}\)，\(n = 1\)，\(2\)，\(3\)，\(\cdots\).

\begin{enumerate}
\item 证明 \(\left\{b_{n}\right\}\) 为等比数列；
\item 如果数列 \(\{b_n\}\) 前 \(3\) 项的和等于 \(\frac{7}{24}\)，求数列 \(\{a_n\}\) 的首项 \(a_1\) 和公差 \(d\).
\end{enumerate}
\end{problem}

\begin{solution}
设等差数列 \(\{a_n\}\) 的首项为 \(a_1>0\)，公差为 \(d\)，则 \(a_n=a_1+(n-1)d\). 由 \(\lg a_1\)、\(\lg a_2\)、\(\lg a_4\) 成等差数列，得 \(2\lg a_2=\lg a_1+\lg a_4\). 由于各项为正，可以利用对数运算将其化为 \(a_2^2=a_1a_4\). 代入 \(a_2=a_1+d\)、\(a_4=a_1+3d\)，得 \((a_1+d)^2=a_1(a_1+3d)\)，即 \(d^2-a_1d=0\)，所以 \(d=0\) 或 \(d=a_1\).

\begin{enumerate}
\item 当 \(d=0\) 时，\(a_n=a_1\)，从而 \(b_n=\frac1{a_{2^n}}=\frac1{a_1}\)，所以 \(\frac{b_{n+1}}{b_n}=1\)，\(\{b_n\}\) 是公比为 \(1\) 的等比数列.

当 \(d=a_1\) 时，\(a_n=a_1+(n-1)a_1=na_1\)，所以 \(a_{2^n}=2^na_1\)，从而 \(b_n=\frac1{2^na_1}\). 因此 \(\frac{b_{n+1}}{b_n}=\frac12\)，\(\{b_n\}\) 是公比为 \(\frac12\) 的等比数列. 两种情形下结论均成立.
\item 当 \(d=0\) 时，\(b_1=b_2=b_3=\frac1{a_1}\)，由题设 \(\frac3{a_1}=\frac7{24}\)，解得 \(a_1=\frac{72}{7}\)，此时 \(d=0\).

当 \(d=a_1\) 时，\(b_1=\frac1{2a_1}\)、\(b_2=\frac1{4a_1}\)、\(b_3=\frac1{8a_1}\)，所以 \(\frac1{2a_1}+\frac1{4a_1}+\frac1{8a_1}=\frac7{8a_1}=\frac7{24}\)，解得 \(a_1=3\)，此时 \(d=3\). 题目没有排除公差为零的情形，故两组解均应保留：\((a_1,d)=\left(\frac{72}{7},0\right)\) 或 \((a_1,d)=(3,3)\).
\end{enumerate}
\end{solution}

\begin{problem}
如图，四棱锥 \(P - ABCD\) 中，底面 \(ABCD\) 为矩形，\(PD \perp\) 底面 \(ABCD\)，\(AD = PD\)，\(E\)，\(F\) 分别为 \(CD\)，\(PB\) 的中点.

\begin{enumerate}
\item 求证：\(EF \perp\) 平面 \(PAB\)；
\item 设 \(AB = \sqrt{2}BC\)，求 \(AC\) 与平面 \(AEF\) 所成的角的大小.
\end{enumerate}

\bitmapfigure[width=0.42\linewidth]{img/2005/national_paper_2_liberal/q20_fig1.png}
\end{problem}

\begin{solution}
设 \(BC=AD=PD=b>0\)，则由 \(AB=\sqrt2BC\) 得 \(AB=\sqrt2b\). 建立空间直角坐标系，取 \(A=(0,0,0)\)，\(B=(\sqrt2b,0,0)\)，\(D=(0,b,0)\)，\(C=(\sqrt2b,b,0)\)，\(P=(0,b,b)\). 由中点坐标公式，\(E=\left(\frac{\sqrt2b}{2},b,0\right)\)，\(F=\left(\frac{\sqrt2b}{2},\frac b2,\frac b2\right)\).

\begin{enumerate}
\item 有 \(\overrightarrow{EF}=\left(0,-\frac b2,\frac b2\right)\)，\(\overrightarrow{AB}=(\sqrt2b,0,0)\)，\(\overrightarrow{AP}=(0,b,b)\). 于是 \(\overrightarrow{EF}\cdot\overrightarrow{AB}=0\)，且 \(\overrightarrow{EF}\cdot\overrightarrow{AP}=0\). 直线 \(AB\) 与 \(AP\) 相交于 \(A\)，并且都在平面 \(PAB\) 内，因此同时垂直于这两条相交直线的直线 \(EF\) 垂直于平面 \(PAB\).
\item 平面 \(AEF\) 内两条相交直线 \(AE\)、\(AF\) 的方向向量分别为 \(\overrightarrow{AE}=\left(\frac{\sqrt2b}{2},b,0\right)\)、\(\overrightarrow{AF}=\left(\frac{\sqrt2b}{2},\frac b2,\frac b2\right)\). 取它们的叉积的非零倍数 \(\bs{n}=(2,-\sqrt2,-\sqrt2)\) 作为平面 \(AEF\) 的法向量. 又 \(\overrightarrow{AC}=C-A=b(\sqrt2,1,0)\)，所以 \(\left|\overrightarrow{AC}\cdot\bs{n}\right|=\sqrt2b\)，\(\left|\overrightarrow{AC}\right|=\sqrt3b\)，\(\left|\bs{n}\right|=\sqrt{4+2+2}=2\sqrt2\). 设 \(AC\) 与平面 \(AEF\) 所成的角为 \(\theta\). 直线与平面所成角的正弦等于直线方向向量与平面法向量夹角余弦的绝对值，因此
\(\sin\theta=\frac{\left|\overrightarrow{AC}\cdot\bs{n}\right|}{\left|\overrightarrow{AC}\right|\left|\bs{n}\right|}=\frac{\sqrt2b}{\sqrt3b\cdot2\sqrt2}=\frac{1}{2\sqrt3}=\frac{\sqrt3}{6}\). 所以 \(\theta=\arcsin\frac{\sqrt3}{6}\).
\end{enumerate}
\end{solution}

\begin{problem}
设 \(a\) 为实数，函数 \(f(x) = x^3 - x^2 - x + a\).

\begin{enumerate}
\item 求 \(f(x)\) 的极值；
\item 当 \(a\) 在什么范围内取值时，曲线 \(y = f(x)\) 与 \(x\) 轴仅有一个交点.
\end{enumerate}
\end{problem}

\begin{solution}
\begin{enumerate}
\item 求导得 \(f'(x)=3x^{2}-2x-1=(3x+1)(x-1)\). 当 \(x<-\frac13\) 时，两个因式均为负，故 \(f'(x)>0\)；当 \(-\frac13<x<1\) 时，两个因式异号，故 \(f'(x)<0\)；当 \(x>1\) 时，两个因式均为正，故 \(f'(x)>0\). 因此函数在 \(x=-\frac13\) 处由增变减，取得极大值，在 \(x=1\) 处由减变增，取得极小值. 代入得 \(f\left(-\frac13\right)=-\frac1{27}-\frac19+\frac13+a=a+\frac5{27}\)，以及 \(f(1)=1-1-1+a=a-1\). 所以极大值为 \(a+\frac5{27}\)，极小值为 \(a-1\).
\item 因为 \(f(x)\) 是首项系数为正的三次函数，所以 \(x\to-\infty\) 时 \(f(x)\to-\infty\)，\(x\to+\infty\) 时 \(f(x)\to+\infty\). 若极大值小于零，即 \(a+\frac5{27}<0\)，则函数在两个驻点处都小于零，左侧递增段和中间递减段没有零点，只有右侧递增段有一个零点，故有且仅有一个交点；这给出 \(a<-\frac5{27}\). 若极小值大于零，即 \(a-1>0\)，则函数在两个驻点处都大于零，只有左侧递增段在趋于 \(-\infty\) 时穿过零点，故也有且仅有一个交点；这给出 \(a>1\).

若 \(-\frac5{27}<a<1\)，则极大值为正、极小值为负，函数在三个单调区间内各有一个零点，共有三个交点. 当 \(a=-\frac5{27}\) 时，极大值为零，且右侧递增段还另有一个零点，共有两个不同交点；当 \(a=1\) 时，极小值为零，且左侧递增段还另有一个零点，也共有两个不同交点. 因此所求范围为 \(a<-\frac5{27}\) 或 \(a>1\).
\end{enumerate}
\end{solution}

\begin{problem}
\(P\)，\(Q\)，\(M\)，\(N\) 四点都在椭圆 \(x^{2} + \frac{y^{2}}{2} = 1\) 上，\(F\) 为椭圆在 \(y\) 轴正半轴上的焦点，已知 \(\overrightarrow{PF}\) 与 \(\overrightarrow{FQ}\) 共线，\(\overrightarrow{MF}\) 与 \(\overrightarrow{FN}\) 共线，且 \(\overrightarrow{PF} \cdot \overrightarrow{MF} = 0\). 求四边形 \(PMQN\) 的面积的最小值和最大值.
\end{problem}

\begin{solution}
椭圆可写为 \(\frac{x^{2}}{1^{2}}+\frac{y^{2}}{(\sqrt2)^{2}}=1\)，所以长半轴为 \(\sqrt2\)，短半轴为 \(1\)，焦距的一半为 \(\sqrt{2-1}=1\)，且焦点为 \(F=(0,1)\). 由共线条件，\(P\)，\(F\)，\(Q\) 在同一条过 \(F\) 的弦上，\(M\)，\(F\)，\(N\) 在另一条过 \(F\) 的弦上. 因为 \(F\) 在椭圆内部，所以它位于两条弦的内部；这两条弦正是四边形 \(PMQN\) 的两条对角线. 又 \(\overrightarrow{PF}\cdot\overrightarrow{MF}=0\)，故两条对角线互相垂直.

设第一条弦的单位方向向量为 \((\cos\theta,\sin\theta)\)，则其上任一点可表示为 \((t\cos\theta,1+t\sin\theta)\). 代入椭圆方程得
\(\left(\cos^{2}\theta+\frac12\sin^{2}\theta\right)t^{2}+\sin\theta\,t-\frac12=0\).
记该方程两根为 \(t_{1}\)、\(t_{2}\)，则弦长为根差的绝对值. 由二次方程根差公式，且 \(\cos^{2}\theta+\frac12\sin^{2}\theta=\frac{1+\cos^{2}\theta}{2}\)，有
\(L_{1}=|t_{1}-t_{2}|=\frac{\sqrt{\sin^{2}\theta+2\left(\cos^{2}\theta+\frac12\sin^{2}\theta\right)}}{\cos^{2}\theta+\frac12\sin^{2}\theta}=\frac{2\sqrt2}{1+\cos^{2}\theta}\).

与第一条弦垂直的单位方向向量可取 \((-\sin\theta,\cos\theta)\). 同样计算可得第二条弦长 \(L_{2}=\frac{2\sqrt2}{1+\sin^{2}\theta}\). 两条对角线垂直，所以四边形面积为
\(S=\frac12L_{1}L_{2}=\frac{4}{(1+\cos^{2}\theta)(1+\sin^{2}\theta)}=\frac{4}{2+\sin^{2}\theta\cos^{2}\theta}\).

由 \(0\leqslant\sin^{2}\theta\cos^{2}\theta\leqslant\frac14\)，可得 \(\frac{16}{9}\leqslant S\leqslant2\). 当 \(|\sin\theta|=|\cos\theta|\) 时，\(\sin^{2}\theta\cos^{2}\theta=\frac14\)，面积取最小值 \(\frac{16}{9}\)；当 \(\sin\theta\cos\theta=0\) 时，面积取最大值 \(2\). 因此四边形 \(PMQN\) 的面积最小值为 \(\frac{16}{9}\)，最大值为 \(2\).
\end{solution}
